% -*- coding: utf-8 -*-
% !TEX program = xelatex
\documentclass[plain,noanswer,medmath]{../../template/jnuexam}
\usepackage{../../template/kysx}

\begin{document}


\examtitle{1997}{1}


\makepart{填空题}{本题共5分，每小题3分，满分15分}


\begin{problem}
$\lim_{x \to 0} \frac{3\sin x + x^2\cos \frac{1}{x}}{(1 + \cos x)\ln(1 + x)} =$ \fillin{}．
\end{problem}


\begin{problem}
设幂级数$\sum_{n = 0}^{\infty}a_n x^n$的收敛半径为$3$,
则幂级数$\sum_{n = 1}^{\infty}na_n (x - 1)^{n + 1}$的收敛区间为 \fillin{}．
\end{problem}


\begin{problem}
对数螺线$\rho = \e^{\theta}$在点$(\rho ,\theta) = \left(\e^{\frac{\pi}{2}},\tfrac{\pi}{2} \right)$
处的切线的直角坐标方程为 \fillin{}．
\end{problem}


\begin{problem}
设$A = \begin{pmatrix}
  1 &  2 & -2 \\
  4 &  t & 3 \\
  3 & -1 & 1
\end{pmatrix}$,$B$为三阶非零矩阵，且$AB = O$，则$t =$ \fillin{}．
\end{problem}


\begin{problem}
袋中有$50$个乒乓球，其中$20$个是黄球，$30$个是白球．
今有两人依次随机地从袋中各取一球，取后不放回，则第二个人取得黄球的概率是 \fillin{}．
\end{problem}


\makepart{选择题}{本题共5小题，每小题3分，满分15分}


\begin{problem}
二元函数$f(x,y) = \begin{cases}
 \frac{xy}{x^2 + y^2}, & (x,y) \ne(0,0), \\
 0, & (x,y) = (0,0)
\end{cases}$在点$(0,0)$处\pickout{}
\begin{abcd}
\item 连续，偏导数存在．                
\item 连续，偏导数不存在．
\item 不连续，偏导数存在．              
\item 不连续，偏导数不存在．
\end{abcd}
\end{problem}


\begin{problem}
设在区间$[a,b]$上$f(x) > 0$, $f'(x) < 0$, $f''(x) > 0$，令
$$S_1 = \int_a^b f(x)\dx,\quad S_2 = f(b)(b - a),\quad S_3 = \frac{1}{2}[f(a) + f(b)](b - a),$$
则\pickout{}
\begin{abcd}
\item $S_1 < S_2 < S_3$．                     
\item $S_2 < S_1 < S_3$．
\item $S_3 < S_1 < S_2$．                     
\item $S_2 < S_3 < S_1$．
\end{abcd}
\end{problem}


\begin{problem}
$\text{设}F(x) = \int_x^{x + 2\pi} \e^{\sin t} \sin t\dt$，则$F(x)$\pickout{}
\begin{abcd}
\item 为正常数．    
\item 为负常数．    
\item 恒为零．    
\item 不为常数．
\end{abcd}
\end{problem}


\begin{problem}
设$\alpha_1 =\begin{pmatrix}
 a_1 \\
 a_2 \\
 a_3 \\
\end{pmatrix}$, $\alpha_2 = \begin{pmatrix}
 b_1 \\
 b_2 \\
 b_3 \\
\end{pmatrix}$, $\alpha_3 = \begin{pmatrix}
 c_1 \\
 c_2 \\
 c_3 \\
\end{pmatrix}$，则三条直线
$$a_1 x + b_1 y + c_1 = 0,\quad a_2 x + b_2 y + c_2 = 0,\quad a_3 x + b_3 y + c_3 = 0$$
（其中$a_i^2+ b_i^2\ne 0,i = 1,2,3$）交于一点的充要条件是\pickout{}
\begin{abcd}
\item $\alpha_1,\alpha_2,\alpha_3$线性相关．
\item $\alpha_1,\alpha_2,\alpha_3$线性无关．
\item 秩$r(\alpha_1,\alpha_2,\alpha_3) =$秩$r(\alpha_1,\alpha_2)$．
\item $\alpha_1,\alpha_2,\alpha_3$线性相关，$\alpha_1,\alpha_2$线性无关．
\end{abcd}
\end{problem}


\begin{problem}
设两个相互独立的随机变量$X$和$Y$的方差分别为$4$和$2$，则随机变量$3X - 2Y$的方差是\pickout{}
\begin{abcd}
\item $8$．        
\item $16$．        
\item $28$．        
\item $44$．
\end{abcd}
\end{problem}


\makepart{计算题}{本题共3小题，每小题5分，满分15分}


\begin{problem}
计算$I = \iiint_{\Omega}(x^2 + y^2) \d V$，其中$\Omega$为平面曲线$\begin{cases}
  y^2 = 2z, \\
  x = 0
\end{cases}$绕$z$轴旋转一周形成的曲面与平面$z = 8$所围成的区域．
\end{problem}


\begin{problem}
计算曲线积分$\oint_C (z - y)\dx + (x - z)\dy + (x - y)\dz$，其中$C$是曲线$\begin{cases}
  x^2 + y^2 = 1, \\
  x - y + z = 2,
\end{cases}$从$z$轴正向往$z$轴负向看，$C$的方向是顺时针的．
\end{problem}


\begin{problem}
在某一人群中推广新技术是通过其中已掌握新技术的人进行的．
设该人群的总人数为$N$，在$t = 0$时刻已掌握新技术的人数为$x_0$，
在任意时刻$t$已掌握新技术的人数为$x(t)$（将$x(t)$视为连续可微变量），
其变化率与已掌握新技术人数和未掌握新技术人数之积成正比，比例常数$k > 0$，求$x(t)$．
\end{problem}


\makepart{计算题}{本题共2小题，第(1)小题6分，第(2)小题7分，满分13分}


\begin{problem}
设直线$L:\begin{cases}
  x + y + b = 0, \\
  x + ay - z - 3 = 0
\end{cases}$在平面$\Pi$上，且平面$\Pi$与曲面$z = x^2 + y^2$相切于点$(1, - 2,5)$，求$a,b$之值．
\end{problem}


\begin{problem}
设函数$f(u)$具有二阶连续导数，而$z = f(\e^x\sin y)$满足方程
$\frac{\pd^2 z}{\pd x^2} + \frac{\pd^2 z}{\pd y^2} = \e^{2x} z$，求$f(u)$．
\end{problem}


\makepart{}{本题满分6分}

\begin{problem}
设$f(x)$连续，$\varphi (x) = \int_0^1f(xt) \dt$，且$\lim_{x \to 0} \frac{f(x)}{x} = A$（$A$为常数），求$\varphi'(x)$并讨论$\varphi'(x)$在$x = 0$处的连续性．
\end{problem}


\makepart{}{本题满分8分}

\begin{problem}
设$a_1 = 2,a_{n + 1} = \frac{1}{2}\left(a_n + \frac{1}{a_n} \right),n = 1,2, \cdots$，证明：\par
\begin{enumerate*}
\item $\lim_{n \to \infty} a_n$存在；
\item 级数$\sum_{n = 1}^{\infty}\left(\frac{a_n}{a_{n + 1}} - 1 \right)$收敛．
\end{enumerate*}
\end{problem}


\makepart{解答题}{本题共2小题，第1小题5分，第2小题6分，满分11分}


\begin{problem}
设$B$是秩为$2$的$5 \times 4$矩阵，
$$\alpha_1 = (1,1,2,3)^T,\alpha_2 = (- 1,1,4, - 1)^T,\alpha_3 = (5, - 1, - 8,9)^T$$
是齐次线性方程组$Bx = 0$的解向量，求$Bx = 0$的解空间的一个标准正交基．
\end{problem}


\begin{problem}
已知$\xi = \begin{pmatrix}
  1 \\
  1 \\
  -1
\end{pmatrix}$是矩阵$A = \begin{pmatrix}
   2 & -1 & 2 \\
   5 &  a & 3 \\
  -1 &  b & -2
\end{pmatrix}$的一个特征向量．
\begin{enumerate}
\item 试确定参数$a,b$及特征向量$\xi$所对应的特征值；
\item 问$A$能否相似于对角阵？说明理由．
\end{enumerate}
\end{problem}


\makepart{}{本题满分5分}

\begin{problem}
设$A$是$n$阶可逆方阵，将$A$的第$i$行和第$j$行对换后得到的矩阵记为$B$．\par
\begin{enumerate*}
\item 证明$B$可逆；
\item 求$AB^{-1}$．
\end{enumerate*}
\end{problem}


\makepart{}{本题满分7分}

\begin{problem}
从学校乘汽车到火车站的途中有$3$个交通岗，假设在各个交通岗遇到红灯的事件是相互独立的，
并且概率都是$\frac{2}{5}$．设$X$为途中遇到红灯的次数，求随机变量$X$的分布律、分布函数和数学期望．
\end{problem}


\makepart{}{本题满分5分}

\begin{problem}
设总体$X$的概率密度为
$$f(x) = \begin{cases}
 (\theta + 1)x^{\theta}, & 0 < x < 1, \\
 0, & \text{其它}, \\
\end{cases}$$
其中$\theta > - 1$是未知参数．$x_1,x_2, \cdots ,x_n$是来自总体$X$的一个容量为$n$的简单随机样本，
分别用矩估计法和最大似然估计法求$\theta$的估计量．
\end{problem}


\end{document}
